\section{The pruning lemma}\label{sec:pruning}

\begin{idea}[What this section does]
Given a left family $\Rf$ and a right family $\Cf$ that may conflict, we
delete some members so that no conflicts remain. The pruning lemma bounds
the number of deletions by $\abs H$ plus the number of conflict points that
end up uncovered. The proof has two parts.
(1)~A purely combinatorial reduction: it suffices to show that a certain
bipartite graph $\Gamma$ has no matching larger than $\abs Z+\abs H$
(\cref{sec:gamma}).
(2)~A linear-algebra proof of this matching bound
(\cref{sec:strategy,sec:matrices,sec:rank,sec:generic,sec:kernel}).
Nothing in this section uses orientation. The relation $\to$ on $X$ is
arbitrary.
\end{idea}

\subsection{Statement and examples}

Recall from \cref{sec:pairs} that a left member $r$ covers $\Ff\{r\}$, a right
member $c$ covers $\Fs\{c\}$, and $Z$ and $H$ are the conflict points and
opposite corners of the conflicts between $\Rf$ and $\Cf$
(\cref{def:ZH}).

\begin{lemma}[Pruning lemma; Lemma~3.1 of the original]\label{lem:pruning}
Let $X$ be a finite set with an arbitrary relation $\to$, and let
$\Rf,\Cf\subseteq X\times X$. There are $\Rf'\subseteq\Rf$ and $\Cf'\subseteq\Cf$ such that
\begin{enumerate}[label=\textup{(P\arabic*)}]
\item\label{P1} $\Rf'$ and $\Cf'$ are compatible;
\item\label{P2} every member of $\Rf\cup\Cf$ that conflicts with no member of the
other family is kept;
\item\label{P3} the number of deleted members satisfies
\begin{equation}\label{eq:pruning-bound}
 \abs{\Rf\setminus\Rf'}+\abs{\Cf\setminus\Cf'}\ \le\ \abs H+\bigl\lvert Z\setminus(\Ff\Rf'\cup\Fs\Cf')\bigr\rvert .
\end{equation}
\end{enumerate}
\end{lemma}

The last term counts the conflict points that no kept member covers. The
two families are always counted separately, even when a pair belongs to both.

\begin{example}[One conflict]\label{ex:one-conflict}
If $\Rf=\{r\}$ and $\Cf=\{c\}$ conflict, then $Z=\{z\}$ and $H=\{h\}$ as in
\cref{fig:coord}. Deleting $r$ costs $1\le\abs H$.
\end{example}

\begin{example}[Deleting more can be cheaper]\label{ex:K22}
Let $X=\{0,1,2,3\}$ with the four arrows $0\to1$, $0\to2$, $3\to1$, $3\to2$, and
\[
 \Rf=\{(0,1),(0,2)\},\qquad \Cf=\{(1,3),(2,3)\}
\]
(\cref{fig:K22grid}). A left member $(0,j)$ and a right member $(i,3)$
conflict because $0\to i$ and $3\to j$. So all four pairs conflict, and
\[
 Z=\{1,2\}\times\{1,2\},\qquad H=\{(0,3)\}.
\]
Compatible subfamilies must have $\Rf'=\varnothing$ or $\Cf'=\varnothing$.
Deleting one whole family costs $2$, but then every point of $Z$ is still
covered by the other family: $(i,3)$ covers $(i,1),(i,2)$. So the right side
of~\eqref{eq:pruning-bound} is only $\abs H=1$. The bound fails for the
\emph{cheapest} pruning! It holds for deleting three members ($3\le1+2$) or
all four ($4\le1+4$). The construction in the proof deletes all four.
\end{example}

\begin{figure}[ht]
\centering
\begin{tikzpicture}[scale=1.0]
 \begin{scope}
  \node[vtx] (x0) at (0,0) {$0$};
  \node[vtx] (x1) at (1.2,0.8) {$1$};
  \node[vtx] (x2) at (1.2,-0.8) {$2$};
  \node[vtx] (x3) at (2.4,0) {$3$};
  \draw[arc] (x0)--(x1); \draw[arc] (x0)--(x2);
  \draw[arc] (x3)--(x1); \draw[arc] (x3)--(x2);
 \end{scope}
 \begin{scope}[xshift=4.3cm,yshift=1.6cm]
  \foreach \k in {0,1,2,3} {\node[font=\small] at (\k,0.55) {$\k$}; \node[font=\small] at (-0.6,-\k) {$\k$};}
  \foreach \i in {0,1,2,3} \foreach \j in {0,1,2,3} \node[gridpt] at (\j,-\i) {};
  \foreach \i in {1,2} \foreach \j in {1,2} \node[Zpt] at (\j,-\i) {};
  \node[Hpt] at (3,0) {};
  \node[Lpt] (r1) at (1,0) {}; \node[Lpt] (r2) at (2,0) {};
  \node[Rpt] (c1) at (3,-1) {}; \node[Rpt] (c2) at (3,-2) {};
  \draw[famL,-{Stealth[length=1.8mm]},thick] (r1)--(1,-0.8);
  \draw[famL,-{Stealth[length=1.8mm]},thick] (r2)--(2,-0.8);
  \draw[famR,-{Stealth[length=1.8mm]},thick] (c1)--(2.2,-1);
  \draw[famR,-{Stealth[length=1.8mm]},thick] (c2)--(2.2,-2);
 \end{scope}
 \begin{scope}[xshift=8.7cm,yshift=0.9cm,font=\small]
  \node[Lpt] at (0,0) {}; \node[anchor=west] at (0.3,0) {$\Rf$ (left family)};
  \node[Rpt] at (0,-0.6) {}; \node[anchor=west] at (0.3,-0.6) {$\Cf$ (right family)};
  \node[Zpt] at (0,-1.2) {}; \node[anchor=west] at (0.3,-1.2) {$Z$ (conflict points)};
  \node[Hpt] at (0,-1.8) {}; \node[anchor=west] at (0.3,-1.8) {$H$ (opposite corners)};
 \end{scope}
\end{tikzpicture}
\caption{\Cref{ex:K22}. The left members (blue, row $0$) cover the green
points below them. The right members (orange, column $3$) cover the green
points to their left. Each of the four conflicts has the same opposite
corner $(0,3)$. Removing all conflicts cheaply is impossible relative to
$\abs H=1$. The lemma pays for the deletions with the green points that are
left uncovered.}
\label{fig:K22grid}
\end{figure}

\Cref{ex:K22} shows that \ref{P3} is \emph{not} an upper bound on the
minimum number of deletions. It is an exchange rate. Every deletion beyond
$\abs H$ must be paid for by a conflict point that becomes uncovered. In the
application (\cref{sec:contradiction}) uncovered points count in the
objective exactly like members, so such trades are free.

\begin{example}[Deleting less is necessary]\label{ex:star}
The opposite extreme. Let
$X=\{p_1,\dots,p_n\}\sqcup\{j_1,\dots,j_n\}\sqcup\{i_0,s_0\}$ with the
arrows $p_k\to i_0$ and $s_0\to j_m$ for all $k,m$, and
\[
 \Rf=\{(p_k,j_m):1\le k,m\le n\},\qquad\Cf=\{c\},\quad c=(i_0,s_0).
\]
Every $(p_k,j_m)$ conflicts with $c$, because $p_k\to i_0$ and $s_0\to j_m$
(\cref{fig:star}). So all $n^2+1$ members are in conflicts, while
\[
 Z=\{(i_0,j_m):m\le n\},\qquad H=\{(p_k,s_0):k\le n\},\qquad \abs Z+\abs H=2n.
\]
Deleting every member that is in a conflict costs $n^2+1$, far above the
budget $\abs H+\abs Z=2n$ even when all of $Z$ is left uncovered. Deleting
$c$ alone costs $1\le\abs H$ and keeps every conflict point covered.
Together with \cref{ex:K22}, where everything must go, this shows that a
pruning within the budget has to be chosen with some care: the pruning
lemma is a statement about matchings, not about a greedy rule. (In the
graph $\Gamma$ of the next subsection, $\{c\}\cup Z_R$ is a vertex cover of
size $n+1=\nu(\Gamma)$.)
\end{example}

\begin{figure}[ht]
\centering
\begin{tikzpicture}[scale=1.0]
 \begin{scope}[yshift=-0.9cm]
  \node[vtx] (p1) at (0,0.7) {$p_1$}; \node[vtx] (p2) at (0,-0.7) {$p_2$};
  \node[vtx] (i0) at (1.5,0) {$i_0$};
  \node[vtx] (s0) at (3.2,0) {$s_0$};
  \node[vtx] (j1) at (4.7,0.7) {$j_1$}; \node[vtx] (j2) at (4.7,-0.7) {$j_2$};
  \draw[arc] (p1)--(i0); \draw[arc] (p2)--(i0); \draw[arc] (s0)--(j1); \draw[arc] (s0)--(j2);
 \end{scope}
 \begin{scope}[xshift=8.0cm,yshift=0.1cm]
  \foreach \x/\l in {0/j_1,1/j_2,2/s_0} \node[font=\small] at (\x,0.6) {$\l$};
  \foreach \y/\l in {0/p_1,-1/p_2,-2/i_0} \node[font=\small] at (-0.7,\y) {$\l$};
  \foreach \x in {0,1,2} \foreach \y in {0,-1,-2} \node[gridpt] at (\x,\y) {};
  \node[Lpt] (r11) at (0,0) {}; \node[Lpt] (r12) at (1,0) {}; \node[Lpt] (r21) at (0,-1) {}; \node[Lpt] (r22) at (1,-1) {};
  \node[Rpt] (c) at (2,-2) {};
  \node[Zpt] at (0,-2) {}; \node[Zpt] at (1,-2) {};
  \node[Hpt] at (2,0) {}; \node[Hpt] at (2,-1) {};
  \draw[famL,-{Stealth[length=1.6mm]},thick] (r21)--(0,-1.75);
  \draw[famL,-{Stealth[length=1.6mm]},thick] (r22)--(1,-1.75);
  \draw[famR,-{Stealth[length=1.6mm]},thick] (c)--(1.25,-2);
 \end{scope}
\end{tikzpicture}
\caption{\Cref{ex:star} with $n=2$; only the rows $p_1,p_2,i_0$ and the
columns $j_1,j_2,s_0$ of the grid are drawn. The four left members (blue)
all conflict with the single right member $c$ (orange). The conflict points
(green) lie in row $i_0$ and the opposite corners (purple) in column $s_0$.
Deleting $c$ removes every conflict at cost $1$; deleting the left members
instead would cost $n^2=4$ against $\abs H=2$.}
\label{fig:star}
\end{figure}

\subsection{Reduction to a matching bound}\label{sec:gamma}

\begin{definition}[The augmented graph]\label{def:gamma}
Let $Z_L$ and $Z_R$ be two disjoint copies of $Z$; write $z_L,z_R$ for the
copies of $z\in Z$. The bipartite graph $\Gamma$ has sides
\[
 \Lp=\Rf\sqcup Z_L,\qquad \Rp=\Cf\sqcup Z_R,
\]
and three kinds of edges:
\begin{itemize}
\item[(E1)] $r\,c$ for every conflicting $r\in\Rf$, $c\in\Cf$ (\emph{conflict edges});
\item[(E2)] $r\,z_R$ whenever $z\in Z$ and $r\in\Rf$ covers $z$;
\item[(E3)] $z_L\,c$ whenever $z\in Z$ and $c\in\Cf$ covers $z$.
\end{itemize}
\end{definition}

Every conflict $(r,c)$ with conflict point $z$ produces the path
$z_L - c - r - z_R$ in $\Gamma$, since both $r$ and $c$ cover $z$
(\cref{fig:gamma}). Different conflicts may share vertices and edges. We
write $\nu(\Gamma)$ for the maximum size of a matching.

\begin{figure}[ht]
\centering
\begin{tikzpicture}[
  ov/.style={circle,draw,inner sep=1pt,minimum size=7mm,font=\footnotesize},
  cv/.style={rectangle,draw,inner sep=1pt,minimum size=6mm,font=\footnotesize}]
 \begin{scope}
  \node[font=\footnotesize] at (1,1.2) {$\Lp$}; \node[font=\footnotesize] at (5,1.2) {$\Rp$};
  \node[ov,fill=famL!30] (r) at (1,0.4) {$r$};
  \node[cv,fill=zcol!25] (zl) at (1,-1.2) {$z_L$};
  \node[ov,fill=famR!30] (c) at (5,0.4) {$c$};
  \node[cv,fill=zcol!25] (zr) at (5,-1.2) {$z_R$};
  \draw[thick] (r) -- node[above,font=\footnotesize] {E1} (c);
  \draw[thick,famL] (r) -- node[pos=0.8,above right=-1pt,font=\footnotesize,black] {E2} (zr);
  \draw[thick,famR] (zl) -- node[pos=0.2,below right=-1pt,font=\footnotesize,black] {E3} (c);
  \node[font=\small] at (3,-2.1) {(a) one conflict};
 \end{scope}
 \begin{scope}[xshift=8.6cm,yshift=1.9cm]
  \node[font=\footnotesize] at (0,0.75) {$\Lp$}; \node[font=\footnotesize] at (4,0.75) {$\Rp$};
  \node[ov,fill=famL!30,ultra thick,draw=red!70!black] (r1) at (0,0) {$01$};
  \node[ov,fill=famL!30,ultra thick,draw=red!70!black] (r2) at (0,-0.9) {$02$};
  \node[ov,fill=famR!30,ultra thick,draw=red!70!black] (c1) at (4,0) {$13$};
  \node[ov,fill=famR!30,ultra thick,draw=red!70!black] (c2) at (4,-0.9) {$23$};
  \foreach \z/\y in {11/-2.0,12/-2.75,21/-3.5,22/-4.25} {
   \node[cv,fill=zcol!25] (L\z) at (0,\y) {$\z$};
   \node[cv,fill=zcol!25] (R\z) at (4,\y) {$\z$};
  }
  \foreach \a in {r1,r2} \foreach \b in {c1,c2} \draw (\a)--(\b);
  \draw[famL] (r1)--(R11); \draw[famL] (r1)--(R21);
  \draw[famL] (r2)--(R12); \draw[famL] (r2)--(R22);
  \draw[famR] (L11)--(c1); \draw[famR] (L12)--(c1);
  \draw[famR] (L21)--(c2); \draw[famR] (L22)--(c2);
  \node[font=\footnotesize,text=black!60] at (-0.85,-3.1) {$Z_L$};
  \node[font=\footnotesize,text=black!60] at (4.85,-3.1) {$Z_R$};
  \node[font=\small,align=center] at (2,-5.15) {(b) \cref{ex:K22}; red: a minimum\\ vertex cover ($4$ vertices)};
 \end{scope}
\end{tikzpicture}
\caption{The augmented graph $\Gamma$. (a) A single conflict $(r,c)$ with
conflict point $z$ forces a path of three edges. (b) The graph for
\cref{ex:K22} (pairs written $ij$). Here $\nu(\Gamma)=4\le\abs Z+\abs H=5$. The
minimum vertex cover consists of the four original members, so the proof
deletes all of them.}
\label{fig:gamma}
\end{figure}

\begin{claim}\label{cl:isolated}
A member of $\Rf\cup\Cf$ that is in no conflict is an isolated vertex of $\Gamma$.
\end{claim}

\begin{proof}
Let $r=(p,j)\in\Rf$ be in no conflict. It has no E1 edges. Suppose it had an E2
edge to $z_R$, with $z=(i,j)\in Z$ and $p\to i$. Since $z\in Z$, some conflict has
conflict point $z$, so there is $(i,s)\in\Cf$ with $s\to j$. But then $r$ and
$(i,s)$ conflict, a contradiction. The argument for $\Cf$ is symmetric.
\end{proof}

\begin{claim}[Covers give prunings]\label{cl:cover-to-pruning}
Let $K$ be a vertex cover of $\Gamma$ (every edge has an endpoint in $K$) that
contains no isolated vertex and has $\abs K\le\abs Z+\abs H$. Then
$\Rf'=\Rf\setminus K$ and $\Cf'=\Cf\setminus K$ satisfy \ref{P1}--\ref{P3}.
\end{claim}

\begin{proof}
\ref{P1}: every conflict edge has an endpoint in $K$, which is deleted.
\ref{P2}: by \cref{cl:isolated} such members are isolated, hence not in $K$.
\ref{P3}: let $Z_{\mathrm{cov}}=Z\cap(\Ff\Rf'\cup\Fs\Cf')$. If $z\in Z_{\mathrm{cov}}$ is
covered by a kept $r\in\Rf'$, then the edge $r\,z_R$ forces $z_R\in K$. If it
is covered by a kept $c\in\Cf'$, then $z_L\in K$. Either way, $K$ contains at
least one copy of each $z\in Z_{\mathrm{cov}}$. So
\[
 \abs{\Rf\setminus\Rf'}+\abs{\Cf\setminus\Cf'}=\abs{K\cap(\Rf\sqcup\Cf)}
 \le\abs K-\abs{Z_{\mathrm{cov}}}\le\abs H+\abs Z-\abs{Z_{\mathrm{cov}}},
\]
and $\abs Z-\abs{Z_{\mathrm{cov}}}$ is the last term of~\eqref{eq:pruning-bound}.
\end{proof}

\begin{claim}[K\H{o}nig]\label{cl:konig}
Every bipartite graph has a vertex cover $K$ with $\abs K=\nu$ that contains
no vertex left unmatched by some fixed maximum matching. In particular $K$
contains no isolated vertex.
\end{claim}

\begin{proof}
This is K\H{o}nig's theorem~\cite{konig1931}. We recall the construction to
check the extra property. Fix a maximum matching $\Mm$. Starting from all
unmatched left vertices, explore \emph{alternating paths}: go from left to
right along non-matching edges and from right to left along matching
edges. Let $W_L,W_R$ be the reached vertices, and put
$K=(\Lp\setminus W_L)\cup W_R$.
No unmatched right vertex is reached, since otherwise we would have an
augmenting path. A matched right vertex that is reached leads to its mate.
A matched left vertex can only be reached through its mate. So the two
ends of a matching edge are both reached or both unreached, and $K$ contains
exactly one end of each matching edge. Hence $\abs K=\abs\Mm$. Unmatched left
vertices are reached and so are not in $K$, and unmatched right vertices are
not reached. Finally, $K$ is a cover. An edge from a reached left vertex to
an unreached right vertex cannot be a non-matching edge (it would have been
followed) or a matching edge (whose ends have the same status).
\end{proof}

By \cref{cl:cover-to-pruning,cl:konig}, the pruning lemma follows from:

\begin{proposition}[Matching bound]\label{prop:matching-bound}
$\nu(\Gamma)\le\abs Z+\abs H$.
\end{proposition}

The rest of this section proves \cref{prop:matching-bound}. If there are no
conflicts, $\Gamma$ has no edges and there is nothing to prove, so we
assume there is at least one conflict.

\subsection{Strategy: three kinds of matching edges}\label{sec:strategy}

Let $\Gamma_L$ be the subgraph of E2 edges (between $\Rf$ and $Z_R$) and
$\Gamma_N$ the subgraph of E3 edges (between $Z_L$ and $\Cf$). Fix a maximum
matching $\Mm$ of $\Gamma$ that has, among all maximum matchings, the fewest
conflict edges. Split it as
\[
 \Mm=\underbrace{\Mm_1}_{k\text{ conflict edges}}\ \sqcup\ \underbrace{\alpha}_{\ell\text{ edges of }\Gamma_L}\ \sqcup\ \underbrace{\beta}_{t\text{ edges of }\Gamma_N},
 \qquad \nu(\Gamma)=k+\ell+t,
\]
and let $R_\alpha\subseteq\Rf$, $C_\beta\subseteq\Cf$ be the original members matched
by $\alpha$, respectively $\beta$ (\cref{fig:decomposition}). Let $R_1\subseteq\Rf$ and $C_1\subseteq\Cf$ be
the left and right ends of the conflict edges in $\Mm_1$. Then
$\abs{R_1}=\abs{C_1}=k$, and $R_1\cap R_\alpha=\varnothing$, $C_1\cap C_\beta=\varnothing$,
since a vertex is matched only once.

\begin{figure}[ht]
\centering
\begin{tikzpicture}[font=\small]
 \draw[rounded corners,fill=famL!12] (0,0) rectangle (2.2,3.2); \node at (1.1,3.5) {$\Rf$};
 \draw[rounded corners,fill=zcol!12] (0,-3.2) rectangle (2.2,-0.4); \node at (1.1,-3.55) {$Z_L$};
 \draw[rounded corners,fill=famR!12] (6,0) rectangle (8.2,3.2); \node at (7.1,3.5) {$\Cf$};
 \draw[rounded corners,fill=zcol!12] (6,-3.2) rectangle (8.2,-0.4); \node at (7.1,-3.55) {$Z_R$};
 \node at (-1.0,1.6) {$\Lp$}; \node at (9.2,1.6) {$\Rp$};
 % R1 and R_alpha
 \draw[fill=white] (0.3,2.0) rectangle (1.9,2.9); \node at (1.1,2.45) {$R_1$};
 \draw[fill=white] (0.3,0.3) rectangle (1.9,1.5); \node at (1.1,0.9) {$R_\alpha$};
 \draw[fill=white] (6.3,2.0) rectangle (7.9,2.9); \node at (7.1,2.45) {$C_1$};
 \draw[fill=white] (6.3,0.3) rectangle (7.9,1.5); \node at (7.1,0.9) {$C_\beta$};
 \draw[line width=2.4pt] (1.9,2.45) -- node[above] {$\Mm_1$: $k$ conflict edges} (6.3,2.45);
 \draw[line width=2.4pt,famL] (1.9,0.9) -- node[pos=0.28,below left=-1pt,black] {$\alpha$: $\ell$} (6.0,-1.8);
 \draw[line width=2.4pt,famR] (2.2,-1.8) -- node[pos=0.72,below right=-1pt,black] {$\beta$: $t$} (6.3,0.9);
\end{tikzpicture}
\caption{The three kinds of edges in the maximum matching $\Mm$: conflict
edges (black), edges of $\Gamma_L$ from $\Rf$ to $Z_R$ (blue) and edges of
$\Gamma_N$ from $Z_L$ to $\Cf$ (orange). The linear algebra handles $\ell$ and $t$
by the ranks of two matrices $L,N$, and $k$ by vectors of $\ker L$ that
survive under $B$.}
\label{fig:decomposition}
\end{figure}

\begin{claim}\label{cl:alpha-max}
$\alpha$ is a maximum matching of $\Gamma_L$, and $\beta$ is a maximum matching of $\Gamma_N$.
\end{claim}

\begin{proof}
If $\alpha$ were not maximum in $\Gamma_L$, there would be an $\alpha$-augmenting
path $\pi$ in $\Gamma_L$, from an $\alpha$-unmatched $r\in\Rf$ to an $\alpha$-unmatched
$z_R\in Z_R$. The vertex $z_R$ is unmatched in $\Mm$ too, because all its
neighbors are in $\Rf$, so any $\Mm$-edge at $z_R$ would belong to $\alpha$.
The interior vertices of $\pi$ are matched by $\alpha\subseteq\Mm$. If $r$ is
unmatched in $\Mm$, then $\pi$ is $\Mm$-augmenting, contradicting maximality.
Otherwise $r$ is matched in $\Mm$ by a conflict edge. Removing that edge and
augmenting along $\pi$ gives a maximum matching with fewer conflict edges,
contradicting the choice of $\Mm$. The argument for $\beta$ is symmetric.
\end{proof}

\paragraph{The plan.}
We will put real numbers on the arrows of the relation and build four
matrices $L,N,A,B$ (\cref{sec:matrices}). We will show:
\begin{itemize}
\item the identity $NL=AB$ makes
$\RR^{\Rf}\to\RR^Z\oplus\RR^H\to\RR^{\Cf}$ a \emph{complex}, so the ranks of
its two maps add up to at most $\abs Z+\abs H$ (\cref{sec:rank});
\item $\rank L=\ell$ and $\rank N=t$, using Edmonds' principle that a generic
matrix with prescribed zero pattern has rank equal to the maximum matching
size (\cref{sec:generic});
\item the $k$ conflict edges of $\Mm_1$ add $k$ dimensions to the first
map, through vectors of $\ker L$ that $B$ keeps independent; this uses a
kernel lemma (\cref{sec:kernel}).
\end{itemize}
Together: $\nu(\Gamma)=k+\ell+t\le\abs Z+\abs H$.

\subsection{Four matrices and two factorizations}\label{sec:matrices}

For every arrow $x\to y$ of the relation, introduce two variables $a_{xy}$ and
$b_{xy}$, all distinct. When $x\not\to y$ we set $a_{xy}=b_{xy}=0$. For an
assignment of real values we define four matrices, with rows and columns
indexed by the sets shown:
\[
\begin{array}{ll}
 L\in\RR^{Z\times\Rf}, & L_{(i,j),(p,j)}=a_{pi},\\[2pt]
 N\in\RR^{\Cf\times Z}, & N_{(i,s),(i,j)}=b_{sj},
\end{array}\qquad
\begin{array}{ll}
 B\in\RR^{H\times\Rf}, & B_{(p,s),(p,j)}=b_{sj},\\[2pt]
 A\in\RR^{\Cf\times H}, & A_{(i,s),(p,s)}=a_{pi},
\end{array}
\]
and all entries not of the indicated form (for example
$L_{(i,j),(p,j')}$ with $j\ne j'$) are $0$. In words, the $a$-matrices $L$ and
$A$ move the first coordinate forward along $p\to i$, and the $b$-matrices
$N$ and $B$ move the second coordinate backward along $s\to j$.

\begin{figure}[ht]
\centering
\begin{tikzpicture}
 \begin{scope}
  \node[Lpt,label={[famL!80!black]above left:$r=(p,j)$}] (r) at (0,0) {};
  \node[Hpt,label={[hcol!80!black]above right:$h=(p,s)$}] (h) at (3,0) {};
  \node[Zpt,label={[zcol!70!black]below left:$z=(i,j)$}] (z) at (0,-2.4) {};
  \node[Rpt,label={[famR!80!black]below right:$c=(i,s)$}] (c) at (3,-2.4) {};
  \draw[-{Stealth[length=2.4mm]},thick] (r) -- node[left] {$L:\ a_{pi}$} (z);
  \draw[-{Stealth[length=2.4mm]},thick] (z) -- node[below] {$N:\ b_{sj}$} (c);
  \draw[-{Stealth[length=2.4mm]},thick] (r) -- node[above] {$B:\ b_{sj}$} (h);
  \draw[-{Stealth[length=2.4mm]},thick] (h) -- node[right] {$A:\ a_{pi}$} (c);
 \end{scope}
 \begin{scope}[xshift=8.3cm,yshift=-1.2cm]
  \node (RR) at (0,1.2) {$\RR^{\Rf}$};
  \node (RH) at (2.6,1.2) {$\RR^{H}$};
  \node (RZ) at (0,-1.2) {$\RR^{Z}$};
  \node (RC) at (2.6,-1.2) {$\RR^{\Cf}$};
  \draw[-{Stealth},thick] (RR) -- node[left] {$L$} (RZ);
  \draw[-{Stealth},thick] (RZ) -- node[below] {$N$} (RC);
  \draw[-{Stealth},thick] (RR) -- node[above] {$B$} (RH);
  \draw[-{Stealth},thick] (RH) -- node[right] {$A$} (RC);
  \node at (1.3,0) {$NL=AB$};
 \end{scope}
\end{tikzpicture}
\caption{The two routes from a left member $r$ to a right member $c$. Via the
conflict point $z$ (first change the first coordinate, then the second), or
via the opposite corner $h$ (the other order). For a conflicting pair both
routes give the entry $a_{pi}b_{sj}$, and for a non-conflicting pair both
give $0$. This is the identity $NL=AB$. The corners sit at their true
grid positions (rows $p,i$; columns $j,s$), and the diagram on the right
has the same shape. Arrows show the direction of the
linear maps. $L$ and $A$ follow the arc $p\to i$, while $N$ and $B$ run
\emph{against} the arc $s\to j$, unlike the set operator $\Fs$ in
\cref{fig:coord}.}
\label{fig:square}
\end{figure}

\begin{claim}\label{cl:supports}
The variable entries of $L$ (those not identically zero) sit exactly on the
edges of $\Gamma_L$, and those of $N$ exactly on the edges of $\Gamma_N$.
\end{claim}

\begin{proof}
$L_{z,r}$ with $z=(i,j)$, $r=(p,j)$ is the variable $a_{pi}$ if $p\to i$, i.e.\
if $r$ covers $z$, and is $0$ otherwise. Similarly $N_{c,z}$ with $c=(i,s)$,
$z=(i,j)$ is $b_{sj}$ exactly when $s\to j$, i.e.\ when $c$ covers $z$.
\end{proof}

\begin{claim}[Two factorizations]\label{cl:commute}
For every assignment, $NL=AB$. Both products have entry $a_{pi}b_{sj}$ at
position $(c,r)$ if $r=(p,j)$ and $c=(i,s)$ conflict, and $0$ otherwise.
\end{claim}

\begin{proof}
$(NL)_{c,r}=\sum_{z\in Z}N_{c,z}L_{z,r}$. A nonzero term needs $z$ in row $i$
(for $N$) and in column $j$ (for $L$), so $z=(i,j)$, and then the term is
$b_{sj}a_{pi}$. If $r,c$ conflict, then $(i,j)\in Z$ and the entry is
$a_{pi}b_{sj}$. If not, then $p\not\to i$ or $s\not\to j$, and the product is $0$.
Likewise $(AB)_{c,r}=\sum_{h\in H}A_{c,h}B_{h,r}$ has at most one term,
$h=(p,s)$, with value $a_{pi}b_{sj}$. If $r,c$ conflict, then $(p,s)\in H$. If
not, the value is $0$.
\end{proof}

The matrix $NL=AB$ might be called the \emph{conflict matrix}: it is
supported on the conflict edges. The first factorization passes through
$\RR^Z$ and the second through $\RR^H$. The whole proof plays these two
against each other.

\begin{remark}[The four matrices are two commuting operators]\label{rem:kronecker}
There is a structural reason for \cref{cl:commute}. Let
$\mathsf a,\mathsf b\in\RR^{X\times X}$ be the matrices of the variables,
$(\mathsf a)_{xy}=a_{xy}$ and $(\mathsf b)_{xy}=b_{xy}$ (zero off the relation),
and consider the two operators on $\RR^{X\times X}$, with coordinates
indexed by grid points,
\[
 (\mathcal A x)_{(i,j)}=\sum_{p}a_{pi}\,x_{(p,j)},\qquad
 (\mathcal B x)_{(i,s)}=\sum_{j}b_{sj}\,x_{(i,j)} .
\]
In Kronecker notation $\mathcal A=\mathsf a^{\T}\otimes I$ and
$\mathcal B=I\otimes\mathsf b$: $\mathcal A$ moves the first coordinate
forward along the relation and $\mathcal B$ moves the second coordinate
backward. They are the linear versions of $\Ff$ and of $\Fs$ for the
reversed relation, and they \emph{commute},
$\mathcal A\mathcal B=\mathcal B\mathcal A=\mathsf a^{\T}\otimes\mathsf b$,
simply because they act on different coordinates. This is the linear
shadow of $\Ff\Fs=\Fs\Ff$. The four matrices are restrictions of these
two: $L$ and $A$ are $\mathcal A$ with rows and columns restricted to
$Z\times\Rf$ and $\Cf\times H$, while $N$ and $B$ are $\mathcal B$
restricted to $\Cf\times Z$ and $H\times\Rf$. The $(c,r)$ entry of
$\mathcal B\mathcal A$ is a single term, with middle index $z=(i,j)$, and
the $(c,r)$ entry of $\mathcal A\mathcal B$ is a single term with middle
index $h=(p,s)$; both equal $a_{pi}b_{sj}$. When this is nonzero $r$ and
$c$ conflict, so $z\in Z$ and $h\in H$, and restricting the middle index to
$Z$, respectively $H$, loses nothing. That is $NL=AB$. The whole of
\cref{sec:pruning} can thus be read as linear algebra on the grid, with the
same two commuting moves as in \cref{sec:pairs}.
\end{remark}

\subsection{A three-term complex}\label{sec:rank}

The identity $NL=AB$ says that the square in \cref{fig:square} commutes.
Put the two routes side by side: the maps
\[
 \LB:\RR^{\Rf}\to\RR^{Z}\oplus\RR^{H},\qquad
 \NmA:\RR^{Z}\oplus\RR^{H}\to\RR^{\Cf}
\]
compose to $NL-AB=0$. In other words,
\begin{equation}\label{eq:complex}
 \RR^{\Rf}\xrightarrow{\ \LB\ }\RR^{Z}\oplus\RR^{H}\xrightarrow{\ \NmA\ }\RR^{\Cf}
\end{equation}
is a \emph{complex}: the image of the first map lies in the kernel of the
second.

\begin{claim}[Dimension count]\label{cl:rank-inequality}
For every assignment, $\rank\LB+\rank\NA\le\abs Z+\abs H$.
\end{claim}

\begin{proof}
By~\eqref{eq:complex} and rank--nullity for the second map,
$\rank\LB\le\dim\ker\NmA=\abs Z+\abs H-\rank\NmA$, and
$\rank\NmA=\rank\NA$, since changing the sign of a block of columns does
not change the rank.
\end{proof}

So \cref{prop:matching-bound} follows once we find an assignment with
\begin{equation}\label{eq:two-lower-bounds}
 \rank\LB\ \ge\ \ell+k\qquad\text{and}\qquad \rank\NA\ \ge\ t .
\end{equation}
The second inequality only needs $\rank N=t$, as $N$ is a submatrix of
$\NA$. In the first, $\ell$ comes from $\rank L=\ell$, and the remaining
$k$ dimensions come from the conflict edges of $\Mm_1$: each gives a vector
of $\ker L$ that $B$ does not kill, and these images stay independent. The
next two subsections provide the assignment and the vectors.

\begin{remark}
The original preprint counts in two steps, with rank--nullity for $A$ and
Sylvester's inequality $\rank(NL)\ge\rank L+\rank N-\abs Z$, placing the
$k$ extra dimensions in $\ker A$, where $B$ sends $\ker L$ because
$AB=NL$. The complex~\eqref{eq:complex} packages both steps into
\cref{cl:rank-inequality}, and $\ker A$ is no longer needed. The Lean
formalization follows the two-step count.
\end{remark}

\subsection{Choosing the coefficients}\label{sec:generic}

We use three elementary facts about real polynomials in finitely many
variables. (i)~A product of nonzero polynomials is nonzero. (ii)~A nonzero
polynomial is nonzero at some real point (induction on the number of
variables: a nonzero univariate polynomial has finitely many roots).
(iii)~Let $\Phi$ be a square matrix whose entries are $0$ or variables, with
\emph{distinct} variables in distinct positions. If the positions of the
variables contain a perfect matching (a permutation $\sigma$ with all
$\Phi_{x,\sigma(x)}$ variables), then $\det\Phi\ne0$. Indeed, distinct
permutations give products of distinct sets of variables, so the monomial of
$\sigma$ cannot cancel.

\begin{claim}[Generic rank of $L$ and $N$]\label{cl:edmonds}
\begin{enumerate}[label=(\alph*)]
\item For every assignment, $\rank L\le\ell$ and $\rank N\le t$.
\item The square submatrix of $L$ with columns $R_\alpha$ and rows
$\{z:z_R\text{ is matched by }\alpha\}$ has a determinant that is a nonzero
polynomial. The same holds for the submatrix of $N$ with rows $C_\beta$ and
columns $\{z: z_L\text{ is matched by }\beta\}$.
\end{enumerate}
\end{claim}

\begin{proof}
(a) If a square submatrix of $L$ has nonzero determinant, some term of its
Leibniz expansion is nonzero. That term is a bijection between its rows and
columns using only nonzero entries, that is, by \cref{cl:supports}, a
matching in $\Gamma_L$. So $\rank L\le\nu(\Gamma_L)=\ell$ by \cref{cl:alpha-max}.
Likewise for $N$.

(b) Entries of $L$ vanish unless row and column have the same second
coordinate $j$. So after grouping rows and columns by $j$, the submatrix is
block diagonal, and each block is square because $\alpha$ matches $(p,j)$ with
some $(i,j)$. Within the block of a fixed $j$, the entry in row $(i,j)$ and
column $(p,j)$ is $a_{pi}$, so distinct positions carry distinct variables
(the same variable may reappear in \emph{other} blocks, but this is
harmless). The edges of $\alpha$ give a perfect matching of every block. By
(iii) every block determinant is a nonzero polynomial, and by (i) so is
their product. For $N$ group by the first coordinate $i$ instead, where the
entries $b_{sj}$ in a block are again distinct.
\end{proof}

We need one more family of rank conditions, which are used in
\cref{sec:kernel}. For every arrow $p\to i$ let
\[
 J_p=\{j:(p,j)\in\Rf\},\qquad S_i=\{s:(i,s)\in\Cf\},\qquad
 \Phi^{p,i}=\bigl(b_{sj}\bigr)_{s\in S_i,\ j\in J_p},
\]
a matrix whose entry is the variable $b_{sj}$ when $s\to j$ and $0$
otherwise. (In the original preprint this matrix is called $D^{p,i}$.) Its
variables are distinct. Let $\rho^{p,i}$ be the largest rank it attains over
all real assignments. At an assignment attaining it, some
$\rho^{p,i}\times\rho^{p,i}$ minor is nonzero, so that minor is a nonzero
polynomial (the empty minor, equal to $1$, if $\rho^{p,i}=0$).

\begin{claim}[A simultaneous choice]\label{cl:simultaneous}
There is a real assignment of all variables at which
\begin{enumerate}[label=(\alph*)]
\item the columns of $L$ indexed by $R_\alpha$ form a basis of $\im L$, and $\rank L=\ell$;
\item the rows of $N$ indexed by $C_\beta$ form a basis of the row space of $N$, and $\rank N=t$;
\item every $\Phi^{p,i}$ has its maximum possible rank $\rho^{p,i}$.
\end{enumerate}
\end{claim}

\begin{proof}
Multiply the two polynomials of \cref{cl:edmonds}(b) and one maximal nonzero
minor of each $\Phi^{p,i}$. This is a finite product of nonzero polynomials,
hence nonzero by (i), and by (ii) it is nonzero at some real point. There, the
$\ell$ columns $R_\alpha$ of $L$ are linearly independent, and since
$\rank L\le\ell$ they span $\im L$. Similarly for $N$. Each $\Phi^{p,i}$ has a
nonzero minor of size $\rho^{p,i}$, so its rank is (at least, hence exactly)
$\rho^{p,i}$.
\end{proof}

\emph{From now on we fix such an assignment.} Different conditions share
variables (for example, $b_{sj}$ appears in $N$, in $B$ and in several
$\Phi^{p,i}$). The point of multiplying the polynomials is that all
conditions hold \emph{at the same time}.

\subsection{The kernel lemma and the bound on \texorpdfstring{$\ker A$}{ker A}}\label{sec:kernel}

\begin{lemma}[Kernel lemma; Lemma~3.2 of the original]\label{lem:kernel}
Let $\Phi\in\RR^{S\times J}$ have the maximum rank among all real matrices
whose nonzero entries lie in a given set $T\subseteq S\times J$, and suppose
$\Phi$ itself is supported in $T$. If $x\in\ker\Phi$ and $y\in\ker\Phi^\T$, then
$x_jy_s=0$ for every $(s,j)\in T$.
\end{lemma}

\begin{proof}
Suppose $x_jy_s\ne0$ for some $(s,j)\in T$, and let $\Phi'=\Phi+\gamma e_se_j^\T$ with
$\gamma\ne0$. It is still supported in $T$. Now $\Phi'x=\gamma x_je_s$, so
$e_s\in\im\Phi'$. Also, for every $w$, $\Phi w=\Phi'w-\gamma w_je_s\in\im\Phi'$, so
$\im\Phi\subseteq\im\Phi'$. But $e_s\notin\im\Phi$, because $y$ is orthogonal to
$\im\Phi$ and $y_s\ne0$. Hence $\rank\Phi'>\rank\Phi$, a contradiction.
\end{proof}

\begin{example}
Let $T=\{(1,1),(2,1)\}$ in a $2\times2$ matrix (column $2$ is forced to be
zero) and $\Phi=\begin{psmallmatrix}1&0\\1&0\end{psmallmatrix}$, of the maximum rank $1$. Then
$\ker\Phi=\langle e_2\rangle$ and $\ker\Phi^\T=\langle(1,-1)\rangle$. The pairs
$(s,j)\in T$ have $j=1$, and $x_1=0$ for all $x\in\ker\Phi$, as the lemma predicts.
For the non-generic $\Phi=\begin{psmallmatrix}1&1\\1&1\end{psmallmatrix}$ on the full support (maximum
rank $2$) the conclusion fails: $x=y=(1,-1)$.
\end{example}

So at a generic point the left and right kernels ``avoid each other'' on
the support. We apply this to the matrices $\Phi^{p,i}$.

\begin{remark}[What the kernel lemma says combinatorially]\label{rem:DM}
Read $T$ as a bipartite graph between the rows $S$ and the columns $J$, and
call a column \emph{inessential} if some maximum matching of $T$ misses it;
likewise for rows. If $x\in\ker\Phi$ and $x_j\ne0$, then column $j$ of
$\Phi$ is a combination of the other columns, so deleting it does not lower
the rank. The matrix without column $j$ is supported in $T\setminus\{j\}$,
so its rank is at most $\nu(T\setminus\{j\})$ by the Leibniz argument of
\cref{cl:edmonds}(a), while $\rank\Phi=\nu(T)$. Hence
$\nu(T\setminus\{j\})=\nu(T)$: some maximum matching misses $j$. So
$\ker\Phi$ is supported on inessential columns and $\ker\Phi^\T$ on
inessential rows, and \cref{lem:kernel} says: \emph{there is no edge of
$T$ between an inessential row and an inessential column.} This is a
basic fact of matching theory, part of the Dulmage--Mendelsohn
decomposition of a bipartite graph~\cite{DulmageMendelsohn1958,LovaszPlummer1986},
and the perturbation argument above is a quick proof of it. Only the
maximum rank is needed, not that the entries be generic. In
\cref{cl:conflict-kernel} the graph is the relation $s\to j$ between
$S_i$ and $J_p$, and a conflict $(r,c)$ is one of its edges.
\end{remark}

\begin{claim}[Kernels separate conflicts]\label{cl:conflict-kernel}
For every $u\in\ker B$, every $v\in\ker N^\T$, and every conflicting pair
$r\in\Rf$, $c\in\Cf$,
\[
 u_r\,v_c=0 .
\]
\end{claim}

\begin{proof}
Let $r=(p,j)$ and $c=(i,s)$, so $p\to i$ and $s\to j$. Consider
$\Phi=\Phi^{p,i}$ and the restrictions
\[
 x=(u_{(p,j')})_{j'\in J_p}\in\RR^{J_p},\qquad y=(v_{(i,s')})_{s'\in S_i}\in\RR^{S_i}.
\]
\emph{$x\in\ker\Phi$.} Take a row $s'\in S_i$. If no $j'\in J_p$ has $s'\to j'$, the
row of $\Phi$ is zero. Otherwise, pick such $j'$. Then $(p,j')\in\Rf$ and
$(i,s')\in\Cf$ conflict, so $(p,s')\in H$, and the $(p,s')$ coordinate of
$Bu=0$ reads
\[
 0=(Bu)_{(p,s')}=\sum_{j'\in J_p}b_{s'j'}\,u_{(p,j')}=(\Phi x)_{s'}.
\]
\emph{$y\in\ker\Phi^\T$.} Take a column $j'\in J_p$. If no $s'\in S_i$ has $s'\to j'$,
the column is zero. Otherwise such a conflict gives $(i,j')\in Z$, and the
$(i,j')$ coordinate of $N^\T v=0$ reads
\[
 0=(N^\T v)_{(i,j')}=\sum_{s'\in S_i}b_{s'j'}\,v_{(i,s')}=(\Phi^\T y)_{j'}.
\]
By \cref{cl:simultaneous}(c), $\Phi$ has maximum rank on its support, and
$(s,j)$ is in the support since $s\to j$. By \cref{lem:kernel},
$x_jy_s=u_rv_c=0$.
\end{proof}

Note that \cref{lem:kernel} is applied at the fixed assignment. The
perturbation in its proof is only a thought experiment about $\Phi^{p,i}$
and does not change any of the other matrices.

\begin{claim}\label{cl:kerA}
$\rank\LB\ge\ell+k$.
\end{claim}

\begin{proof}
Let $(r,c)\in\Mm_1$ be a conflict edge of the matching. Since $r\notin R_\alpha$,
\cref{cl:simultaneous}(a) lets us write $Le_r$ as a combination of the
columns $Le_a$, $a\in R_\alpha$. That is, there is a vector
\[
 u^r=e_r+\sum_{a\in R_\alpha}\lambda^r_ae_a\in\ker L .
\]
Similarly, since $c\notin C_\beta$, \cref{cl:simultaneous}(b) gives
$v^c=e_c+\sum_{b\in C_\beta}\mu^c_be_b\in\ker N^\T$, with $v^c_c=1$.

We claim that the $\ell+k$ vectors
\[
 \LB e_a=(Le_a,Be_a)\quad(a\in R_\alpha),\qquad
 \LB u^r=(0,Bu^r)\quad(r\in R_1)
\]
are linearly independent. Suppose
$\sum_{a\in R_\alpha}\theta_a\LB e_a+\sum_{r\in R_1}\theta_r\LB u^r=0$.
The $\RR^Z$-components give $\sum_a\theta_aLe_a=0$, and the columns $Le_a$,
$a\in R_\alpha$, are independent by \cref{cl:simultaneous}(a), so every
$\theta_a=0$. The $\RR^H$-components then give $\sum_r\theta_rBu^r=0$, that
is, $u=\sum_{r\in R_1}\theta_ru^r\in\ker B$. For each $r\in R_1$, with
partner $c$ in $\Mm_1$, \cref{cl:conflict-kernel} applied to $u$ and $v^c$
gives $u_r=u_rv^c_c=0$. But $u_r=\theta_r$, because $u^{r'}$ is supported on
$\{r'\}\cup R_\alpha$ and $R_1\cap R_\alpha=\varnothing$. So all $\theta_r=0$.
\end{proof}

\begin{proof}[Proof of \cref{prop:matching-bound} and of \cref{lem:pruning}]
At the assignment fixed after \cref{cl:simultaneous} we have $\rank N=t$,
and $\rank\LB\ge\ell+k$ by \cref{cl:kerA}. By \cref{cl:rank-inequality},
\[
 \nu(\Gamma)=k+\ell+t\ \le\ \rank\LB+\rank\NA\ \le\ \abs Z+\abs H .
\]
\Cref{cl:konig} gives a vertex cover of this size with no isolated vertex,
and \cref{cl:cover-to-pruning} turns it into the required pruning.
\end{proof}

\begin{idea}[Summary of the pruning lemma]
The matching bound is the heart of the matter. Because $\Ff$ and $\Fs$
commute, the conflict matrix factors through $\RR^Z$ and through $\RR^H$,
$NL=AB$, and this makes $\RR^{\Rf}\to\RR^Z\oplus\RR^H\to\RR^{\Cf}$ a
complex: the ranks of its two maps add up to at most $\abs Z+\abs H$. A
maximum matching of $\Gamma$ certifies lower bounds for these ranks:
$\ell$ edges through $L$ and $t$ edges through $N$ (generic rank $=$
matching number), and $k$ conflict edges, each contributing a vector of
$\ker L$ whose image under $B$ is new. The kernel lemma keeps these images
independent.
\end{idea}

\begin{remark}[Towards a combinatorial proof?]\label{rem:combinatorial}
Neither orientation nor linear algebra appears in the \emph{statement} of
\cref{prop:matching-bound}, and one may ask for a proof by alternating
paths. Here is what a maximum matching sees. With the notation of
\cref{sec:strategy}, let $Z_\alpha$ and $Z_\beta$ be the sets of $z\in Z$
whose copies $z_R$, respectively $z_L$, are matched by $\Mm$. Then
\[
 \nu(\Gamma)=k+\abs{Z_\alpha}+\abs{Z_\beta}
 =k+\abs{Z_\alpha\cup Z_\beta}+\abs{Z_\alpha\cap Z_\beta},
\]
so the bound would follow from $k+\abs{Z_\alpha\cap Z_\beta}\le\abs H$.
Each of the $k+\abs{Z_\alpha\cap Z_\beta}$ objects on the left comes with
an opposite corner: a conflict edge $(r,c)$ has its $h$, and a doubly
matched $z$ is matched to some $r_z$ and some $c_z$, which both cover $z$
and therefore conflict, with a corner $h_z$. If these corners were always
distinct we would be done. They need not be, and the inequality can fail
for a maximum matching with the fewest conflict edges. In a computer
search over about $5000$ random instances with small $\Gamma$ (the folder
\texttt{experiments/}), there was always \emph{some} maximum matching
satisfying $k+\abs{Z_\alpha\cap Z_\beta}\le\abs H$. Whether this holds in
general, and whether it has a direct proof, we do not know. The kernel
lemma itself has a combinatorial meaning (\cref{rem:DM}); the place where
linear algebra seems essential is the \emph{sum} $\sum_r\theta_ru^r$ in the
proof of \cref{cl:kerA}, which handles several alternating structures at
once.
\end{remark}
